\chapter{Network models}
\label{chap:redes}

\section{From a situation to a graph}
A network represents places, states, or people as nodes; arcs represent
possible connections. The meaning of the number on an arc depends on the
model: cost or distance in a route, capacity in a maximum-flow problem.
Confusing these quantities changes the problem \cite{ahuja1993network}.
The examples in this chapter are original, and their files are in
\texttt{manual/examples/}. They all use the JSON envelope with \texttt{kind}
and \texttt{model}; WinQSB files are not required.
Run the commands from the repository root.

\section{Shortest path and minimum spanning tree}
SPP seeks a directed path from the source to the destination with minimum cost.
Dijkstra requires nonnegative costs. In QSB, the first node is the source and
the last is the destination. An absent arc is not equivalent to a zero-cost arc.
The file \path{plan_spp.json} has nodes A, B, C, D and the following arcs:
\begin{center}
\begin{tabular}{lrrrrr}\toprule
Arc & A--B & A--C & B--C & B--D & C--D \\
Cost & 2 & 5 & 1 & 6 & 2 \\
\bottomrule\end{tabular}
\end{center}
\begin{lstlisting}[language=bash]
swift run qsb validate-network-json manual/examples/plan_spp.json
swift run qsb solve-network-json manual/examples/plan_spp.json --backend native
swift run qsb solve-network-json manual/examples/plan_mst.json --backend native
\end{lstlisting}
\begin{figure}[htbp]
\centering
\begin{tikzpicture}
\node[qsb node] (A) at (0,0) {A};
\node[qsb node] (B) at (3,1.2) {B};
\node[qsb node] (C) at (3,-1.2) {C};
\node[qsb node] (D) at (6,0) {D};
\draw[qsb selected] (A) -- node[qsb label,above] {2} (B);
\draw[qsb alternative] (A) -- node[qsb label,below] {5} (C);
\draw[qsb selected] (B) -- node[qsb label,right] {1} (C);
\draw[qsb alternative] (B) -- node[qsb label,above] {6} (D);
\draw[qsb selected] (C) -- node[qsb label,below] {2} (D);
\end{tikzpicture}

\caption{Shortest-path network: labels are costs. Thick arrows trace the optimal route; dashed arrows show alternatives.}
\label{fig:shortest-path}
\end{figure}

SPP returns A $\to$ B $\to$ C $\to$ D with cost $2+1+2=5$.
MST connects all nodes without cycles at the lowest total cost; it does not
seek a path between endpoints. The example uses the same connections as
undirected edges, and Kruskal chooses B--C, A--B, C--D, with cost 5.
QSB normalizes asymmetric MST connections to the lower cost for each pair;
a disconnected network has no tree spanning all nodes.

\section{Maximum flow}
With capacities $u_{ij}$, flow must satisfy
$0\leq x_{ij}\leq u_{ij}$ and conservation at intermediate nodes. The
objective maximizes net flow leaving the source. It does not minimize cost.
In \path{plan_maxflow.json}, S--A has capacity 3, S--B 2,
A--B 1, A--T 2, and B--T 3. Although the JSON field is named \texttt{cost},
it represents capacity in MFP.
\begin{lstlisting}[language=bash]
swift run qsb solve-network-json manual/examples/plan_maxflow.json --backend native
\end{lstlisting}
\begin{figure}[htbp]
\centering
\begin{tikzpicture}
\node[qsb node] (S) at (0,0) {S};
\node[qsb node] (A) at (3,1.2) {A};
\node[qsb node] (B) at (3,-1.2) {B};
\node[qsb node] (T) at (6,0) {T};
\draw[qsb arc] (S) -- node[qsb label,above] {3/3} (A);
\draw[qsb arc] (S) -- node[qsb label,below] {2/2} (B);
\draw[qsb arc] (A) -- node[qsb label,right] {1/1} (B);
\draw[qsb arc] (A) -- node[qsb label,above] {2/2} (T);
\draw[qsb arc] (B) -- node[qsb label,below] {3/3} (T);
\draw[densely dotted,line width=1pt] (0.75,-1.7) -- (0.75,1.7);
\end{tikzpicture}

\caption{Maximum flow: each label gives flow/capacity. The dotted line separates the source from the remaining nodes and marks a saturated cut.}
\label{fig:maximum-flow}
\end{figure}

Edmonds--Karp obtains flow 5: 3 and 2 units leave S, and 2 and 3 reach T.
The cut separating S from the other nodes has capacity 5 and proves optimality.
Again, the first and last nodes determine the endpoints.

\section{Traveling salesperson}
TSP seeks a directed cycle that visits every node exactly once and returns
to the starting node. Connecting the nodes is insufficient: there must be
a single tour. The example \path{plan_tsp.json} contains three nodes, with
$c_{AB}=2$, $c_{BC}=3$, $c_{CA}=4$ and reverse costs 4, 5, and 6.
\begin{lstlisting}[language=bash]
swift run qsb solve-network-json manual/examples/plan_tsp.json --backend native
\end{lstlisting}
The tour A $\to$ B $\to$ C $\to$ A costs 9; the reverse tour costs 15.
Held--Karp is exact for the accepted model but exponential; the solver
rejects more than 20 nodes. A model with no feasible Hamiltonian cycle is
not repaired by adding connections. The \texttt{fixtureScale} label expresses
this scale limitation, not a guarantee of fast execution near the maximum.

\section{Assignment and transportation}
Assignment minimizes $\sum_{ij}c_{ij}x_{ij}$, assigning each worker to one
task and using each task at most once. The Hungarian method accepts
rectangular matrices with at least as many tasks as workers.
In \path{plan_assignment.json}, the rows Ana and Luis have costs
$(2,7,4)$ and $(6,1,3)$ for X, Y, and Z.
\begin{lstlisting}[language=bash]
swift run qsb solve-network-json manual/examples/plan_assignment.json --backend native
swift run qsb solve-network-json manual/examples/plan_transport.json --backend native
\end{lstlisting}
Ana--X and Luis--Y cost 3; Z remains unassigned. The model does not include
teams, multiple tasks per person, or additional constraints.

Transportation allows continuous quantities and requires balance:
\[
 \min\sum_{ij}c_{ij}x_{ij},\qquad
 \sum_jx_{ij}=s_i,\quad \sum_ix_{ij}=d_j,\quad x_{ij}\geq0.
\]
The example has supply $(4,6)$, demand $(5,5)$, and cost matrix
$\left(\begin{smallmatrix}1&4\\3&2\end{smallmatrix}\right)$.
QSB ships 4 from O1 to D1, 1 from O2 to D1, and 5 from O2 to D2;
the cost is $4+3+10=17$. Check row and column sums separately.
An imbalance is a validation error in TP. This family does not automatically
create a dummy source or destination.

\section{Minimum-cost flow}
CNF combines transshipment and costs: at each node,
$\sum_jx_{ij}-\sum_jx_{ji}=s_i-d_i$, and it minimizes
$\sum_{ij}c_{ij}x_{ij}$ with nonnegative continuous quantities.
The current model has no arc capacities or lower bounds.
The validator requires finite nonnegative costs in CNF and the graph
variants; assignment permits finite negative costs. Do not model a subsidy
as a negative cost on a CNF arc, because it will be rejected.
\path{plan_netflow.json} supplies 4 at O, requires 4 at D, and allows
O--M at cost 1, M--D at cost 2, and O--D at cost 5.
\begin{lstlisting}[language=bash]
swift run qsb solve-network-json manual/examples/plan_netflow.json --backend native
\end{lstlisting}
The 4 units pass through M and cost 12; direct shipment would cost 20.
CNF preserves original supply and demand and, when they differ, emits a
warning and adds a zero-cost dummy adjustment during solving.
Always review \texttt{balanceAdjustments}: a solution with an adjustment
does not prove that original supply can meet actual demand.

\section{macOS use and backend selection}
Open the example by importing normalized JSON. In Definition,
SPP, MST, MFP, and TSP have a node and arc editor;
AP and TP have cost matrices, with supply and demand in TP; CNF has nodes,
arcs, and balances. Changes are validated before Run. The endpoint order
for SPP/MFP maps to the JSON order; visual positions are not exported.
The solution Diagram highlights the path, tree, tour, or connections used.
\texttt{All connections} also shows the remaining connections. In bipartite
networks, the groups represent workers/tasks or sources/destinations.
Use the detail panel to read exact values and JSON to audit the result;
the thickness of a connection does not replace its numerical quantity.

Native offers deterministic algorithms for all seven variants. Validate
checks the model without solving it. External is available for TP and CNF
through HiGHS; there is no external backend covering all network variants.
Solutions preserve the model, algorithm, assumptions, and exactness. The
practical size of native LPs depends on the formulation; for large
transportation networks, consider HiGHS and check its availability.

\section{Exercises}
\begin{enumerate}
\item Change A--C from 5 to 1 in SPP. Verify the new path and explain why
      the minimum spanning tree may respond differently.
\item Reduce S--A to 1 in MFP. Calculate a cut before solving.
\item Change D2 demand to 6 in TP. Observe validation; then
      deliberately represent one additional unit of supply.
\item In CNF, increase demand to 5. Interpret the dummy adjustment and
      explain which business data are missing before the plan can be accepted.
\end{enumerate}
